\documentclass[11pt]{article}
\usepackage[margin=1in]{geometry}
\usepackage{amsmath}
\usepackage{amssymb}
\usepackage{booktabs}
\usepackage{graphicx}
\usepackage{hyperref}

\title{Time Series Forecasting Project Report\\
\large Global Temporal Retrieval with an Inverted Transformer Backbone}
\author{Group ID: 34 \\ Members: Mert Uyar, Efe Akkus}
\date{}

\begin{document}
\maketitle

\begin{abstract}
We study multivariate hourly forecasting on a panel of 96 related series, where the goal
is to predict a target channel over a 336 hour horizon from a short window of recent
observations and a set of covariates. Our model combines a Global Temporal Retriever (GTR) \cite{cao2026gtr},
which stores a learnable bank of periodic patterns and fuses them with the input through a
small two dimensional convolution, with an inverted Transformer (iTransformer) backbone that
treats each channel as a token. The long horizon is produced autoregressively in 24 hour
blocks so that known future covariates can be used at every step. The dominant period is
selected with an autocorrelation analysis, as suggested for cycle based models, and we compare
cycle lengths of 24, 168, and 672 hours. On the internal validation split the best
configuration reaches a Weighted Absolute Percentage Error (WAPE) of about 30.69 with a daily
cycle. We find that the recurrent 24 hour rollout clearly beats a single direct 336 step head,
that a shorter cycle and a smaller learning rate help most, and that a good learning rate and
gradient clipping matter more for this task than the choice of cycle length. The GitHub repo: 
\href{https://github.com/MertUyar/multivariate_time_series_forecasting}{\texttt{https://github.com/MertUyar/multivariate\_time\_series\_forecasting}}.
\end{abstract}

\section{Introduction}
Forecasting many correlated time series at once is a common industrial problem. In our setting
we are given 96 units, each recorded hourly, and for every unit we observe a target signal
together with a group of covariates such as calendar features, workload and demand estimates,
and several risk indicators. The task is to forecast the target of each unit for the next
336 hours, that is two weeks ahead, using only the recent past and the covariates that are
known for the horizon.

The problem is interesting for three reasons. First, the horizon is long relative to the input
window, so a model has to capture structure that repeats over days and weeks rather than only
short local dynamics. Second, the 96 units share behaviour but differ in level and scale, which
means a single global model has to normalise per unit while still learning shared patterns.
Third, several covariates are given for the future, so a good method should exploit this known
future information instead of ignoring it. These properties make the dataset a useful test bed
for models that mix explicit periodic structure with a learned sequence model.

Our contributions are as follows. We implement a two stage model in PyTorch that first passes
the input through a Global Temporal Retriever \cite{cao2026gtr} and then through a forecasting backbone, and we
support both an inverted Transformer and a decomposition linear model as interchangeable
backbones. Our main goal is to find out whether GTR \cite{cao2026gtr} is a good model for this task since it is 
a proven plug and play model to capture long periodicites. We describe how the periodic pattern bank is retrieved by cycle phase and fused with
the input by a convolution, and how the long horizon is produced by an autoregressive rollout
that reuses known future covariates. We select the modelling period with an autocorrelation
study and run ablations over cycle length, look back length, learning rate, loss function, and
the horizon strategy. Finally we report the validation results and discuss what mattered most.

\section{Related Work}
Transformer models have become a standard tool for sequence modelling since the introduction of
self attention \cite{vaswani2017}. For long horizon forecasting a line of work has adapted the
architecture to reduce cost and to better match the structure of time series, including
Informer \cite{zhou2021informer} and Autoformer \cite{wu2021autoformer}. A different and
simpler view was put forward by DLinear \cite{zeng2023dlinear}, which decomposes a series into
trend and seasonal parts and applies a linear map to each, and which remains a strong and
cheap baseline. PatchTST \cite{nie2023patchtst} showed that treating patches of a single channel
as tokens is very effective. The iTransformer \cite{liu2024itransformer} inverts the usual
layout and treats each variate as a token, so that attention operates across channels rather
than across time. This inversion fits multivariate panels well and is the main backbone we use.

A separate thread models periodicity explicitly. RevIN \cite{kim2022revin} removes per instance
mean and variance before the model and restores them afterwards, which stabilises training when
series drift in level. CycleNet \cite{lin2024cyclenet} argues that many series are dominated by
a small number of periods and proposes to learn a shared bank of periodic components indexed by
the position within the cycle. Our Global Temporal Retriever \cite{cao2026gtr} follows this idea: it keeps a
learnable table of periodic patterns and retrieves the entry that matches the current phase,
which lets the backbone focus on the residual behaviour rather than on re learning the cycle at
every window. Following the CycleNet recommendation, we identify the period with an
autocorrelation analysis instead of fixing it by hand.

\section{Method}
\subsection{Problem setup}
Let a single unit have a history $x_{1:T}\in\mathbb{R}^{T\times N}$, where $T$ is the look back
length and $N$ is the number of channels, with the target channel placed last. We forecast the
target over a horizon $H$. The backbone predicts a block of length $S$ at a time, and the full
horizon $H$ is produced by repeating the block prediction. In our runs $T=96$, $S=24$, and
$H=336$, so $H$ is covered by $14$ successive blocks. The forecasting mode is multivariate to
univariate: all $N$ channels enter the model but only the target is scored.

\subsection{Global Temporal Retriever \cite{cao2026gtr}}
The Global Temporal Retriever (GTR) \cite{cao2026gtr} is applied to the input before the backbone. It stores a
learnable parameter bank $Q\in\mathbb{R}^{C\times N}$, where $C$ is the cycle length and $N$ is
the number of channels. Row $c$ of $Q$ is the pattern that the model expects at phase $c$ of the
cycle, for every channel. Given the phase of the first step of the window, $\phi$, the retrieved
pattern for the window is
\begin{equation}
q_{t} = Q\big[(\phi + t)\bmod C\big], \qquad t=0,\dots,T-1,
\end{equation}
which produces a tensor $q\in\mathbb{R}^{T\times N}$ aligned in phase with the input. A linear
layer along the time axis then refines $q$, giving the model freedom to reshape the retrieved
template rather than copy it exactly.

The refined pattern is fused with the observed input through a small two dimensional
convolution. We stack the input and the retrieved pattern along a new axis of size two, so each
channel is represented by a $2\times T$ image whose first row is the observation and whose
second row is the retrieved cyclic template. A convolution with a single input and output
channel, a kernel of height two and width $2\lfloor P/2\rfloor+1$, and temporal padding
$\lfloor P/2\rfloor$ mixes the two rows across a temporal window of roughly one period $P$:
\begin{equation}
z = \mathrm{Conv2D}\big(\,[\,x \,;\, q\,]\,\big), \qquad
\hat{x} = z + x .
\end{equation}
The kernel height of two lets the convolution combine the current value with the phase matched
global pattern, the kernel width tied to the period lets it borrow information from neighbouring
positions within the cycle, and the residual connection keeps the original signal intact so the
retriever only has to add periodic structure. A dropout of $0.1$ is applied inside GTR \cite{cao2026gtr}. The
output has the same shape as the input and is passed to the backbone.

\subsection{Forecasting backbones}
\paragraph{iTransformer.} Our main backbone is an inverted Transformer
\cite{liu2024itransformer}. Each channel is embedded into a token by projecting its length $T$
series to a $d_{\text{model}}$ vector, so the sequence of tokens has length equal to the number
of channels. Temporal features and a one hot encoding of the unit identity are appended as
extra tokens, which lets attention share information across channels and across units. An
encoder only stack of self attention and feed forward blocks processes the tokens, and a linear
projector maps each channel token to the $S$ step forecast, after which the covariate tokens are
discarded and only the real channels are kept. We use $d_{\text{model}}=512$,
$d_{\text{ff}}=512$, three encoder layers, eight heads, and dropout $0.1$.

\paragraph{DLinear.} For comparison we replace the backbone with DLinear
\cite{zeng2023dlinear}. It decomposes the input into a trend part, obtained with a moving average
of kernel size 25, and a seasonal remainder, then applies a separate linear map from length $T$
to length $S$ to each part and sums the two. DLinear has no attention and few parameters, which
makes it a useful lower complexity reference for the same pipeline.

\subsection{Normalisation and input encoding}
Because the 96 units differ in level and scale, we apply reversible instance normalisation
\cite{kim2022revin} around the backbone. For each window we subtract the per channel temporal
mean and divide by the per channel temporal standard deviation, run GTR \cite{cao2026gtr} and the backbone on the
normalised input, then multiply by the standard deviation and add the mean back to the forecast:
\begin{equation}
\tilde{x} = \frac{x-\mu}{\sqrt{\sigma^{2}+\varepsilon}}, \qquad
\hat{y} = \sqrt{\sigma^{2}+\varepsilon}\;\hat{y}_{\text{norm}} + \mu .
\end{equation}
We tested with it but it didn't improve or even hurt the training, so decided not to use it.

\subsection{Autoregressive rollout}
The backbone predicts $S=24$ hours at a time, while the required horizon is $H=336$ hours. We
produce the horizon recurrently. Starting from the last $T$ real observations we predict the
next 24 hours, write the predicted target back into the target channel of the input, advance
the window by 24 hours, and repeat until the 336 hour horizon is filled. The covariates for the
horizon are known and are read from the validation input at every step, so only the target
channel is fed back. Predicting short blocks and rolling forward keeps each individual
prediction close to the training distribution and lets the model use the known future covariates
at every step, which we found to be far more accurate than training a single head that outputs
all 336 steps at once.

\subsection{Preprocessing}
Following the CycleNet \cite{lin2024cyclenet} recommendation for cycle based models, we do not fix the period by hand.
As a preprocessing step we detrend each unit and compute the autocorrelation function of the
target, then inspect the correlation at candidate lags of one day, one week, and multiples of a
week. The daily lag of 24 hours and the weekly lag of 168 hours show the strongest and most
consistent autocorrelation across units, so we treat the cycle length as a hyperparameter and
test 24, 168, and 672 hours in the experiments. The data was also scaled with StandardScaler from sklearn library.

\subsection{Training}
We train with the Huber loss with $\delta=0.5$, which is less sensitive to large residuals than
the squared error, using the Adam optimiser. Gradients are clipped to a maximum norm of $1.0$
for stability, which we added after observing unstable early training at higher learning rates.
We use early stopping on the internal validation loss with a patience of five epochs, and we fix
the random seed for Python, NumPy, and PyTorch so that runs are reproducible. The dataset is
read once and cached locally, and inference loads a saved checkpoint, so the final forecasting
step does not require internet access.

\section{Experiments}
\subsection{Data and protocol}
The data consist of 96 hourly units with a shared feature schema of 22 covariates and one target
channel, giving 23 channels per window. For each unit we hold out the last 336 hours of the
training file as an internal validation set and train on the remainder, and we report metrics on
this internal split. The official validation input is used only to produce the final 336 hour
forecast through the autoregressive rollout described above. Unless stated otherwise all runs use
a look back of 96 hours, RevIN not used, a 24 hour prediction block, and the hyperparameters in
Table~\ref{tab:hparams}.

\begin{table}[h]
\centering
\begin{tabular}{lr}
\toprule
Hyperparameter & Value \\
\midrule
Look back length $T$ & 96 \\
Prediction block $S$ & 24 \\
Forecast horizon $H$ & 336 \\
Channels $N$ (\texttt{enc\_in}) & 23 \\
Backbone & iTransformer (DLinear for comparison) \\
$d_{\text{model}}$, $d_{\text{ff}}$ & 512, 512 \\
Encoder layers, heads & 3, 8 \\
Dropout (GTR and backbone) & 0.1 \\
Cycle length $C$ & 24, 168, or 672 \\
Convolution period $P$ & 24 \\
Normalisation & No RevIN \\
Loss & Huber ($\delta=0.5$); MSE for comparison \\
Optimiser & Adam \\
Learning rate & $5\times10^{-4}$ default, $1\times10^{-4}$ best \\
Batch size & 16 (32 for DLinear) \\
Max epochs / patience & 30 / 5 \\
Gradient clipping (max norm) & 1.0 \\
Random seed & 2024 \\
\bottomrule
\end{tabular}
\caption{Main hyperparameters. Lower level defaults follow the released configuration.}
\label{tab:hparams}
\end{table}

We evaluate with the Weighted Absolute Percentage Error, together with mean absolute error and
mean squared error, computed on the target channel after de normalisation:
\begin{equation}
\mathrm{WAPE} = \frac{\sum_i |y_i-\hat{y}_i|}{\sum_i y_i}, \qquad
\mathrm{MAE} = \frac{1}{n}\sum_i |y_i-\hat{y}_i|, \qquad
\mathrm{MSE} = \frac{1}{n}\sum_i (y_i-\hat{y}_i)^2 .
\end{equation}
Lower is better for all three. WAPE is our primary metric because it is scale aware across units.

\subsection{Baselines and backbone comparison}
Table~\ref{tab:main} (the errors are taken from the submissions to the leader board) compares 
our model against one reference forecasts and against the DLinear
backbone inside the same pipeline. The naive last value forecast repeats the last observed
target. These give the minimum and a stronger reference respectively.

\begin{table}[h]
\centering
\begin{tabular}{lrrr}
\toprule
Method & WAPE & MAE & MSE \\
\midrule
Naive last value baseline & 61.35 & 5.29 & 48.61 \\
GTR + DLinear backbone & 33.67 & 3.70 & 25.70 \\
GTR + iTransformer backbone (ours) & \textbf{30.69} & \textbf{3.37} & \textbf{22.87} \\
\bottomrule
\end{tabular}
\caption{Validation results on the internal split, target channel, 336 hour horizon by
autoregressive rollout. Our best configuration uses a daily cycle and a learning rate of
$1\times10^{-4}$. Baseline and DLinear numbers are to be filled from your own runs. 
(The errors are taken from the submissions to the leader board.)}
\label{tab:main}
\end{table}

\subsection{Cycle length}
Table~\ref{tab:cycle} (the errors are taken from the validation split from the training data) 
varies the cycle length while keeping everything else fixed. At the default
learning rate the three cycles are close, with the weekly cycle slightly ahead, but the best
result overall is the daily cycle once the learning rate is reduced. This is consistent with the
autocorrelation analysis, where the 24 hour lag is the most stable across units, and it suggests
that most of the recurring structure the model needs is daily.

\begin{table}[h]
\centering
\begin{tabular}{lrrr}
\toprule
Cycle length (hours) & WAPE & MAE & MSE \\
\midrule
24, learning rate $5\times10^{-4}$ & 25.50 & 2.73 & 14.90 \\
168, learning rate $5\times10^{-4}$ & 24.80 & 2.66 & 14.00 \\
672, learning rate $5\times10^{-4}$ & 25.20 & 2.70 & 14.75 \\
24, learning rate $1\times10^{-4}$ & \textbf{24.40} & \textbf{2.61} & \textbf{13.28} \\
\bottomrule
\end{tabular}
\caption{Cycle length ablation with a 96 hour look back and a 24 hour block. Lower is better. 
(The errors are taken from the validation split from the training data.)}
\label{tab:cycle}
\end{table}

\subsection{Horizon strategy and look back length}
Table~\ref{tab:horizon} (the errors are taken from the validation split from the training data) 
isolates two design choices. Producing the horizon with a single direct
336 step head is clearly worse than the autoregressive 24 hour rollout, which supports our choice
to predict short blocks and reuse the known future covariates at each step. Increasing the look
back from 96 to 336 hours did not help and slightly hurt, so we keep the 96 hour window, which
also keeps the comparison across backbones fair.

\begin{table}[h]
\centering
\begin{tabular}{lrrr}
\toprule
Setting & WAPE & MAE & MSE \\
\midrule
Autoregressive rollout, block 24 & \textbf{25.20} & \textbf{2.70} & \textbf{14.75} \\
Direct head, 336 steps at once & 30.50 & 3.26 & 21.13 \\
Look back 336 (rollout) & 26.10 & 2.79 & 15.28 \\
\bottomrule
\end{tabular}
\caption{Horizon strategy and look back length, at cycle length 672 where both variants were
run. The rollout beats the direct head and the shorter look back is at least as good. 
(The errors are taken from the validation split from the training data.)}
\label{tab:horizon}
\end{table}

\subsection{Learning rate, loss, and stability}
Early in the project the training was unstable: with a learning rate of $5\times10^{-4}$ and a
schedule that held the rate high for the first epochs, the validation loss was best at the first
epoch and then rose, so early stopping ended the run before the model recovered. Both the
training and validation loss increased together, which points to a learning rate that is too
large rather than to overfitting. Reducing the learning rate to $1\times10^{-4}$ and adding
gradient clipping removed this behaviour and gave the best result in Table~\ref{tab:cycle} 
(the errors are taken from the validation split from the training data), and a
constant high rate was clearly worse. Switching the loss from Huber to MSE 
changed the result only marginally, so we kept the Huber loss for its robustness to
large residuals.

\section{Conclusion}
We presented a two stage forecasting model that pairs a Global Temporal Retriever \cite{cao2026gtr} with an
inverted Transformer backbone, and fills a 336 hour horizon
by an autoregressive 24 hour rollout that reuses known future covariates. The main lessons are that recurrent 
short block prediction beats a single long head, that a shorter cycle for this dataset and a smaller learning 
rate help most, and that gradient clipping with a moderate learning rate is needed for stable training. 
Limitations include a look back of only 96 hours, a single random
seed per configuration in the reported numbers. On the internal validation split the best configuration 
reaches a WAPE of about 30.69 with a daily cycle. Unfortunately, GTR \cite{cao2026gtr} with iTransformer \cite{liu2024itransformer} 
backbone disappointed us in his prediction performance. We think that the problem lies in the fact that 
GTR \cite{cao2026gtr} is thought to be for datasets with long periodicity which doesn't apply for this 
dataset. GTR \cite{cao2026gtr} head seems to be only adding additional complexity. Future work could test 
for ablation where GTR \cite{cao2026gtr} is removed, sweep the seed for variance estimates, try a longer 
look back with the daily cycle, and use the covariate forecasts more directly inside the model rather than 
only through the rollout.

\section{LLMs}
Claude is used to help write the report. No LLM is used to implement the model.

\section*{Contributions}
Mert Uyar: Coded GTR \cite{cao2026gtr}. Coded data loaders and experiments. Run experiments. Wrote report. \\
Efe Akkus: Coded the backbones. Coded experiments. Run experiments. Wrote report.

\begin{thebibliography}{10}

\bibitem{cao2026gtr}
F. Cao, L. Dai, J. Han, and H. Xiong. Enhancing Multivariate Time Series Forecasting with Global Temporal Retrieval.
\textit{International Conference on Learning Representations}, 2026. arXiv:2602.10847.

\bibitem{vaswani2017}
A. Vaswani, N. Shazeer, N. Parmar, J. Uszkoreit, L. Jones, A. N. Gomez, L. Kaiser, and
I. Polosukhin. Attention Is All You Need. \textit{Advances in Neural Information Processing
Systems}, 2017.

\bibitem{liu2024itransformer}
Y. Liu, T. Hu, H. Zhang, H. Wu, S. Wang, L. Ma, and M. Long. iTransformer: Inverted Transformers
Are Effective for Time Series Forecasting. \textit{International Conference on Learning
Representations}, 2024.

\bibitem{zeng2023dlinear}
A. Zeng, M. Chen, L. Zhang, and Q. Xu. Are Transformers Effective for Time Series Forecasting?
\textit{AAAI Conference on Artificial Intelligence}, 2023.

\bibitem{lin2024cyclenet}
S. Lin, W. Lin, X. Hu, W. Wu, R. Mo, and H. Zhong. CycleNet: Enhancing Time Series Forecasting
through Modeling Periodic Patterns. \textit{Advances in Neural Information Processing Systems},
2024.

\bibitem{kim2022revin}
T. Kim, J. Kim, Y. Tae, C. Park, J.-H. Choi, and J. Choo. Reversible Instance Normalization for
Accurate Time Series Forecasting against Distribution Shift. \textit{International Conference on
Learning Representations}, 2022.

\bibitem{nie2023patchtst}
Y. Nie, N. H. Nguyen, P. Sinthong, and J. Kalagnanam. A Time Series is Worth 64 Words: Long Term
Forecasting with Transformers. \textit{International Conference on Learning Representations}, 2023.

\bibitem{wu2021autoformer}
H. Wu, J. Xu, J. Wang, and M. Long. Autoformer: Decomposition Transformers with Auto Correlation
for Long Term Series Forecasting. \textit{Advances in Neural Information Processing Systems}, 2021.

\bibitem{zhou2021informer}
H. Zhou, S. Zhang, J. Peng, S. Zhang, J. Li, H. Xiong, and W. Zhang. Informer: Beyond Efficient
Transformer for Long Sequence Time Series Forecasting. \textit{AAAI Conference on Artificial
Intelligence}, 2021.

\end{thebibliography}

\end{document}
